\section{Kapitel 8 -- Fazit und
Ausblick}\label{kapitel-8-fazit-und-ausblick}

\begin{center}\rule{0.5\linewidth}{0.5pt}\end{center}

\subsection{8.1 Zusammenfassung der
Ergebnisse}\label{zusammenfassung-der-ergebnisse}

Die vorliegende Arbeit hat ein hybrides Redaktionssystem für die
KI-gestützte Liveticker-Erstellung im Profifußball konzipiert,
implementiert und evaluiert. Das System entstand in direkter Kooperation
mit der \textbf{Stackwork GmbH} im IT-Bereich von Eintracht Frankfurt
und exportiert produzierte Inhalte über API-Schnittstellen in die
bestehende \textbf{Stackwork Demo App}. Es adressiert die drei in
Kapitel 1.1 identifizierten Problemdimensionen --- operativer Zeitdruck,
Mehrsprachigkeit und White-Label-Bedarf --- durch eine dreischichtige
Architektur aus Datenbeschaffung (n8n), Anwendungslogik (FastAPI,
PostgreSQL) und Präsentation (React, TypeScript).

\subsubsection{8.1.1 Technische Ergebnisse}\label{technische-ergebnisse}

Die technische Evaluation (Kapitel 6.2) belegt eine produktionsnahe
Codebasis mit vollständiger Testpyramide, hoher Typsicherheit und
lückenloser Erfüllung aller 23 definierten Anforderungen (vgl. Kapitel
6.6). Die bewusst ausgeklammerte Authentifizierungsschicht ist als
Systembeschränkung dokumentiert (vgl. Kapitel 6.7).

\subsubsection{8.1.2 KI-Textgenerierung}\label{ki-textgenerierung}

Die Multi-Provider-Architektur mit Prioritätskette gewährleistet eine
hohe Verfügbarkeit der Textgenerierung. Die qualitative Bewertung
(Abschnitt 6.3.6) zeigt durchgängig gute Ergebnisse über alle drei
Qualitätsdimensionen; die stärkste Einschränkung liegt in der
Stil-Inkonsistenz des neutralen Profils. Die quantitativen Einzelwerte
und Latenzmetriken sind in Kapitel 6.3--6.5 vollständig dokumentiert.

\subsubsection{8.1.3 Systemarchitektur und
Designentscheide}\label{systemarchitektur-und-designentscheide}

Das System realisiert die in Kapitel 4 konzipierte dreischichtige
Architektur vollständig (Datenschicht via n8n-ETL, Anwendungsschicht via
FastAPI/PostgreSQL, Präsentationsschicht via React/TypeScript). Die
\textbf{White-Label-Architektur} und die \textbf{drei Betriebsmodi}
bilden die Kernbeiträge des Systemdesigns (vgl. Kap. 4.3.3, 5.3.3).

\begin{center}\rule{0.5\linewidth}{0.5pt}\end{center}

\subsection{8.2 Beantwortung der
Forschungsfrage}\label{beantwortung-der-forschungsfrage}

Die in Kapitel 1.2 formulierte Forschungsfrage lautet:

\begin{quote}
\emph{Inwiefern reduziert ein hybrides KI-gestütztes Redaktionssystem
die Time-to-Publish bei der Liveticker-Erstellung im Profifußball im
Vergleich zur rein manuellen Erstellung, ohne die journalistische
Qualität hinsichtlich Korrektheit, Tonalität und Verständlichkeit zu
beeinträchtigen?}
\end{quote}

Die Beantwortung erfolgt entlang der beiden in der Forschungsfrage
angelegten Dimensionen.

\subsubsection{8.2.1 Zeitliche Dimension: Reduktion der
Time-to-Publish}\label{zeitliche-dimension-reduktion-der-time-to-publish}

Im \texttt{auto}-Modus (vollautonom) entfällt die manuelle
Texterstellung vollständig --- die geschätzte TTP beträgt \textbf{≈
3,4--5,9 s} (erwarteter Median bis Worst Case, vgl. Abschnitt 6.4.2). Im
\texttt{coop}-Modus (hybrid) addiert sich die redaktionelle Prüfzeit:
Ein einfacher Freigabe-Klick erfordert ca. 5--10 s, ein bearbeiteter
Entwurf ca. 15--30 s --- die geschätzte TTP liegt damit bei \textbf{≈
15--30 s}. Im \texttt{manual}-Modus (Status quo) muss der Redakteur den
gesamten Text selbst verfassen; auf Basis der in Kapitel 2.1
beschriebenen Produktionsbedingungen wird die typische
Texterstellungszeit unter Livebedingungen auf \textbf{30--120 s}
geschätzt.

Da kein kontrolliertes Spielexperiment in allen drei Modi durchgeführt
werden konnte, basiert der Vergleich auf gemessenen LLM-Latenzen und der
implementierten Systemarchitektur (Abschnitt 6.4.2). Ein
Cliff's-Delta-Test auf real gemessenen TTP-Paaren (auto vs.~manual) wäre
mit einem zukünftigen Live-Spieltest durchführbar; die TTP-Metrik und
die Bulk-Evaluationsinfrastruktur (Abschnitte 6.3.1--6.3.2) sind dafür
vorbereitet.

Die strukturellen Daten bestätigen die Hypothese, dass ein hybrides
System die Publikationslatenz signifikant reduziert: Die geschätzte
TTP-Reduktion im \texttt{auto}-Modus gegenüber \texttt{manual} beträgt
eine Größenordnung (Faktor 5--35× je nach Event-Typ und
Redakteurserfahrung).

\subsubsection{8.2.2 Qualitative Dimension: Journalistische
Qualität}\label{qualitative-dimension-journalistische-qualituxe4t}

Die journalistische Qualität wurde entlang der drei in der
Forschungsfrage definierten Kriterien operationalisiert:

\begin{enumerate}
\def\labelenumi{\arabic{enumi}.}
\item
  \textbf{Korrektheit}: Die generierten Texte erreichen einen Ø-Wert von
  \textbf{4,6 / 5}. In 15 von 16 bewerteten Einträgen (94 \%) wurden
  alle verfügbaren Fakten (Spieler, Team, Minute, Ergebnis) korrekt
  wiedergegeben (Hauptevaluation, N = 16, deutschsprachig). Die
  ergänzende Evaluation mit synthetischem Kontext (N = 9) bestätigt
  diesen Befund und erweitert ihn um bislang fehlende Event-Typen (Rote
  Karte, Phasen-Events). Die Pre-Match-Schutzregel verhindert
  zuverlässig die Erfindung von Live-Spielszenen; eine Wettempfehlung (1
  von 16 Einträgen, 6 \%) stellt eine inhaltliche Halluzination
  geringerer Schwere dar. Die Qualitätsbewertung basiert ausschließlich
  auf deutschsprachigen Ausgaben; die Texte in anderen Sprachen
  (Englisch, Japanisch) wurden nicht separat evaluiert.
\item
  \textbf{Tonalität}: Mit einem Ø-Wert von \textbf{4,1 / 5} erzeugen die
  drei Stilprofile deutlich unterscheidbare Textstile. Der euphorische
  Modus --- primär für vereinsnahe Liveticker wie die
  \texttt{ef\_whitelabel}-Instanz konzipiert --- erzielt die stärksten
  qualitativen Ergebnisse: Die Beispieltexte in Abschnitt 6.3.6 zeigen,
  dass emotionale Stilmittel des Liveticker-Genres (Wiederholungen,
  Ausrufe, szenische Abschlüsse) vom Modell zuverlässig umgesetzt
  werden. Das neutrale Profil zeigt die größte Inkonsistenz (19 \% der
  Einträge zu emotional formuliert), bedingt durch die starken
  euphorischen Few-Shot-Referenzen.
\item
  \textbf{Verständlichkeit}: Mit \textbf{4,3 / 5} erfüllen die
  generierten Texte die linguistischen Anforderungen des
  Liveticker-Genres (Kapitel 2.5) zuverlässig: kurze Satzkonstruktionen,
  Präsenskonstruktionen, idiomatische Ausrufe und konzeptionelle
  Mündlichkeit sind durchgehend vorhanden. Kein Eintrag überschritt die
  genretypische Kürze.
\end{enumerate}

Die qualitative Dimension der Bewertung basiert auf der in Abschnitt
6.3.6 dokumentierten Selbstevaluation, ergänzt durch das
Experteninterview mit einem professionellen Sportredakteur von Eintracht
Frankfurt (vgl. Abschnitt 6.3.6).

\subsubsection{8.2.3 Synthese}\label{synthese}

Die Forschungsfrage verknüpft zwei Dimensionen durch ein „ohne'': eine
TTP-Reduktion soll erreicht werden, \emph{ohne} die journalistische
Qualität zu beeinträchtigen. Aus der Zusammenschau der Ergebnisse lässt
sich eine modus-spezifische Antwort ableiten.

Im \texttt{auto}-Modus ist die TTP-Reduktion maximal --- das System
operiert ohne menschliche Latenz. Zugleich sind die Qualitätsrisiken in
diesem Modus nicht abgefangen: Die 19\%ige Stil-Inkonsistenz im
neutralen Profil und die 6\%ige Rate inhaltlicher Halluzinationen würden
ohne redaktionelle Filterung direkt publiziert. Die Bedingung „ohne
Beeinträchtigung der Qualität'' ist im vollautonomen Betrieb damit nur
dann erfüllt, wenn ausschließlich das euphorische Profil (das
konsistenteste der drei) verwendet wird --- also im vereinsspezifischen
\texttt{ef\_whitelabel}-Kontext.

Im \texttt{coop}-Modus hingegen werden beide Bedingungen der
Forschungsfrage simultan erfüllt: Die KI-Generierung reduziert die
Texterstellungszeit auf 3,4--5,9 s; die redaktionelle Prüfung (5--30 s)
filtert Qualitätsrisiken heraus, bevor ein Eintrag publiziert wird. Die
resultierende TTP liegt bei 15--30 s --- erheblich unter der manuellen
Baseline von 30--120 s und dennoch qualitätsgesichert. Das
Human-in-the-Loop-Design löst damit den klassischen
Speed-Quality-Trade-off nicht durch Kompromiss, sondern durch
Aufgabenteilung: Die Maschine übernimmt die zeitkritische
Textproduktion, der Mensch die publizistische Verantwortung.

Die Antwort auf die Forschungsfrage lautet daher zustimmend, aber mit
einer Qualifikation: \textbf{Ein hybrides KI-gestütztes Redaktionssystem
reduziert die TTP signifikant, ohne die journalistische Qualität zu
beeinträchtigen --- vorausgesetzt, es wird im \texttt{coop}-Modus
betrieben.} Der \texttt{auto}-Modus maximiert die Geschwindigkeit auf
Kosten der Qualitätssicherung und eignet sich für Kontexte mit geringem
Qualitätsrisiko (euphorisches Profil, vereinseigener Kanal); der
\texttt{manual}-Modus bietet volle redaktionelle Kontrolle ohne
Zeitgewinn. Der \texttt{coop}-Modus realisiert den in Kapitel 1.1
formulierten Zielzustand eines Systems, in dem KI-Geschwindigkeit und
menschliche Urteilskraft gemeinsam wirken.

\begin{center}\rule{0.5\linewidth}{0.5pt}\end{center}

\subsection{8.3 Ausblick}\label{ausblick}

\subsubsection{8.3.1 Kurzfristige
Erweiterungen}\label{kurzfristige-erweiterungen}

\textbf{Authentifizierung und Rollensystem}: Die Integration eines JWT-
oder OAuth-2.0-basierten Authentifizierungsframeworks ist die wichtigste
Voraussetzung für einen Mehrbenutzerbetrieb. Ein rollenbasiertes
Zugriffskonzept (Redakteur, Chefredakteur, Administrator) würde die
Freigabe-Workflows im \texttt{coop}-Modus um eine Vier-Augen-Prüfung
erweitern.

\textbf{Server-Sent Events für Ticker-Updates}: Die Ablösung des
5-Sekunden-Pollings durch SSE für den Ticker-Datenfluss würde die
End-to-End-Latenz im \texttt{coop}-Modus weiter reduzieren und die
Serverlast bei vielen gleichzeitigen Nutzern senken.

\textbf{Erweiterte E2E-Tests}: Die Ausweitung der Playwright-Tests auf
den vollständigen Redaktionsworkflow mit laufendem Backend und
Testdatenbank (Spiel auswählen, Events empfangen, Ticker-Einträge
bearbeiten, Freigeben) würde die Regressionssicherheit erhöhen.

\textbf{n8n-Workflow-Tests}: Die 15 n8n-Workflows (vgl. Kapitel 6.7)
sind derzeit ausschließlich manuell verifiziert. Integrationstests gegen
eine Staging-Datenbank würden die Zuverlässigkeit der
Datenversorgungspipeline absichern.

\textbf{Few-Shot-Kontextualisierung in n8n-Workflows}: Die n8n-Workflows
für Halbzeit- und Abpfiff-Zusammenfassungen (Workflow 13) rufen
OpenRouter direkt auf, ohne Stilreferenzen aus der
\texttt{style\_references}-Tabelle einzubinden. Eine Erweiterung dieser
Workflows um einen SQL-Schritt zur Referenzabfrage würde die
stilistische Kohärenz zwischen Backend-generierten und n8n-generierten
Einträgen herstellen.

\textbf{Größere Evaluationsstichprobe}: Die Hauptevaluation umfasst 16
Einträge aus 9 Spielen --- eine für explorative Erkenntnisse
ausreichende, für statistisch belastbare Aussagen jedoch zu kleine
Datenbasis. Eine Folgestudie mit ≥ 100 Einträgen aus verschiedenen
Wettbewerben, Ligen und Spielsituationen würde die Generalisierbarkeit
der Ergebnisse stärken und den Einsatz der implementierten
nicht-parametrischen Verfahren (Cliff's Delta, Bootstrap-CI)
rechtfertigen.

\textbf{Langzeitbetrieb}: Das System wurde nicht über einen längeren
Zeitraum evaluiert. Aspekte wie Drift der Textqualität über die Zeit,
kumulative API-Kosten im Saison-Betrieb und Rate-Limit-Häufigkeit unter
realer Dauerlast sind nicht erfasst. Eine Saisonbegleitung über
mindestens 34 Bundesliga-Spieltage würde hier belastbare Betriebsdaten
liefern.

\subsubsection{8.3.2 Mittelfristige
Erweiterungen}\label{mittelfristige-erweiterungen}

\textbf{Externe Nutzerstudie}: Eine systematische Evaluation mit
professionellen Sportredakteuren, die das System im realen Spielbetrieb
nutzen, würde die externe Validität der Qualitätsbewertung stärken. Das
implementierte Evaluationsframework (TTP, Cliff's Delta, Cohen's Kappa)
ist dafür vorbereitet.

\textbf{Fine-Tuning statt Few-Shot}: Für vereinsspezifische Instanzen
könnte ein Fine-Tuning auf dem Korpus historischer Vereinsticker die
stilistische Konsistenz über den Few-Shot-Ansatz hinaus verbessern. Die
\texttt{style\_references}-Tabelle bildet bereits eine kuratierte
Datenbasis, die als Ausgangspunkt für ein Fine-Tuning-Dataset dienen
könnte.

\textbf{Differenziertes Polling}: Die vorbereitete Infrastruktur in
\texttt{resolvePollingInterval} ermöglicht eine phasenabhängige
Anpassung der Polling-Frequenz (z. B. 30 Sekunden im PreMatch, 5
Sekunden während Live). Dies würde die Serverlast im Ruhezustand
reduzieren.

\textbf{Evaluation der Mehrsprachigkeit}: Eine separate
Qualitätsevaluation der mehrsprachigen Textgenerierung --- insbesondere
für die im Kontext von Eintracht Frankfurt relevanten Sprachen Englisch
und Japanisch --- würde die in Kapitel 2.2 motivierte
Internationalisierungsanforderung empirisch absichern.

\subsubsection{8.3.3 Langfristige
Perspektiven}\label{langfristige-perspektiven}

\textbf{Multimodale Ticker}: Die bestehende Integration von
Medieninhalten (ScorePlay, YouTube, Instagram) könnte zu einem
multimodalen Ticker weiterentwickelt werden, in dem Bild, Video und Text
automatisch zu einem kohärenten Narrativ verwoben werden.
Vision-Language-Modelle könnten dabei Spielszenen aus Bildern
beschreiben und in den Tickertext integrieren.

\textbf{Autonomes Scoring der Textqualität}: Die Ablösung der manuellen
Qualitätsbewertung durch ein trainiertes Scoring-Modell, das
Korrektheit, Tonalität und Verständlichkeit automatisch bewertet, würde
eine kontinuierliche Qualitätsüberwachung im Produktivbetrieb
ermöglichen. Die implementierte
\texttt{aggregate\_quality\_by\_group}-Funktion bildet die
Aggregationslogik hierfür bereits ab.

\textbf{Generalisierung über den Fußball hinaus}: Die
Provider-agnostische LLM-Architektur und das Event-basierte
Generierungsmodell sind nicht fußballspezifisch. Eine Erweiterung auf
andere Sportarten (Handball, Basketball, Eishockey) erforderte primär
die Anpassung der Event-Typ-Mappings, der Phasendefinitionen und der
Kontextbuilder --- die Kernarchitektur bliebe unverändert.

\begin{center}\rule{0.5\linewidth}{0.5pt}\end{center}

\subsection{8.4 Schlusswort}\label{schlusswort}

Die vorliegende Arbeit zeigt, dass ein hybrides KI-gestütztes
Redaktionssystem den operativen Zeitdruck bei der Liveticker-Erstellung
adressieren kann, ohne die redaktionelle Kontrolle aufzugeben. Der
\texttt{coop}-Modus realisiert das in Kapitel 1.1 formulierte Zielbild
eines Systems, in dem „die finale Entscheidungshoheit und publizistische
Verantwortung beim Menschen verbleiben''.

Die technische Reife des Systems (vgl. Kap. 6.2) sowie die Erfüllung
aller definierten Anforderungen dokumentieren die Tragfähigkeit des
Architekturansatzes. Die in Kapitel 7 diskutierten Limitationen
markieren klare Erweiterungspfade, stellen aber die grundsätzliche
Funktionalität nicht in Frage.

Die KI-gestützte Liveticker-Generierung ersetzt den Redakteur nicht ---
sie gibt ihm die Zeit zurück, die er braucht, um seiner eigentlichen
Aufgabe nachzukommen: Journalismus.
